% Short industry-RAG intro paragraph. Drop in near Table~\ref{tab:rag} (results)
% or at the end of subsec:workloads (methodology). Self-contained.
\paragraph{Industry-RAG workloads.}
Beyond the synthetic length/cache ladders, OpenDC-Infer-LC includes two workloads
shaped like the retrieval traffic that dominates enterprise GenAI deployments.
\textbf{RAG-TopK} reproduces vector-database top-$k$ retrieval: each request packs
$64$ independently retrieved passages ($\approx$32K tokens) from \emph{different}
documents, exactly one of which carries the answer while the other passages are
\emph{hard distractors} in the identical format---the same near-duplicate, wrong-chunk
noise a production retriever returns---and the retrieved set differs per query.
\textbf{RAG-Report} reproduces long-form report generation, today's most common
enterprise LLM task: a 32K multi-document context that must be synthesized into a
structured $2048$-token answer reproducing five required facts. These are
\emph{industry-representative} on three axes that toy long-context tests miss.
\emph{(i)~Length and shape:} 32K is the realistic production RAG operating window
(retrieval budgets, not 4--8K toys), and TopK's fragmented many-document prefill
plus per-query-distinct passage sets reproduce the low cross-query cache reuse real
retrievers exhibit---unlike a single shared haystack. \emph{(ii)~Correctness under
distraction:} grading is exact-match against an answer embedded among same-format
distractors, so a system is rewarded only for retrieving the \emph{right} chunk---the
dominant real RAG failure mode, which raw-throughput benchmarks cannot see.
\emph{(iii)~Both cost regimes:} TopK is prefill/retrieval-bound and Report is
decode-bound on long structured output, the two regimes that jointly set the serving
bill of a production RAG stack. Construction details are in
Appendix~\ref{subsec:rag_details}.
